\chapter*{Conclusion}
\addcontentsline{toc}{chapter}{Conclusion}
\label{ch:conclusion}

\section*{Summary of findings}

On an unbalanced panel of 535 large U.S.\ non-financial firms over
fiscal years 2010--2024, the joint and decomposition specifications
of Chapter~\ref{ch:methodology} deliver three findings.
\textbf{H1 is supported}: holding firm fundamentals, industry, and
year fixed, the natural log of the Item~1A word count is positively
associated with post-filing log-volatility at the 30-day horizon
($\hat{\beta} = +0.130$, $t = +7.70$ in the H1+H3 specification).
The volume effect of \textcite{campbell2014} is therefore present in
a panel that is different in size and time period from the
original, and survives within-firm absorption.
\textbf{H2 is supported}: holding length and the same controls
fixed, the composite LM risk-word density of Item~1A is negatively
associated with post-filing log-volatility ($\hat{\beta} = -3.88$,
$t = -2.74$). The specificity effect of \textcite{hope2016},
translated into dictionary terms, is present once the volume
channel is partialled out.
\textbf{H3 is supported}: decomposing the composite RiskDensity
into its two LM sub-components reveals that the entire negative
effect is carried by the LM-Litigious sub-list
($\hat{\beta}^{\text{lit}} = -6.16$, $t = -3.35$), while the
LM-Uncertainty sub-list is statistically indistinguishable from
zero. The pattern survives at every post-filing horizon between 5
and 365~days and on both COVID-era and pre-COVID sub-samples. The
interpretation aligns with the
\textcite{nelson2007,cazier2021,beatty2019} reading of
litigious-list language as standardised cautionary boilerplate
characteristic of mature, legally supervised disclosure
environments.

\section*{Methodological contribution}

The thesis's methodological contribution has two layers. First,
the dictionary-based risk-word density used in much prior work
mixes two opposite-signed channels into one statistic through a
shared denominator. The univariate density coefficient at the
30-day horizon ($-7.78$) is roughly twice the joint estimate
($-3.88$); the gap is exactly the part of the volume channel that
loads onto density through the negative cross-sectional
correlation between length and density. Entering the two
regressors jointly in a single specification recovers two effects
that the dictionary-only ratio hides.

Second, the H2 ``specificity'' effect itself pools two LM word
lists with very different production technologies. The LM-list
decomposition (equation~\ref{eq:decomposition_model}) shows that
the negative coefficient on the composite is driven entirely by
the LM-Litigious sub-list, while the LM-Uncertainty sub-list adds
no explanatory power. To our knowledge no prior panel study of
post-filing realised volatility estimates the LM-Litigious density
of Item~1A as a stand-alone regressor; the closest references
\parencite{loughran2011, kravet2013, campbell2014, lyle2023}
either aggregate the two sub-lists, hand-code their own risk
categories, or focus on year-on-year changes. The decomposition
therefore localises a known specificity effect to a specific LM
sub-list and gives it a concrete interpretation in light of the
boilerplate-disclosure literature
\parencite{nelson2007, cazier2021, beatty2019}.

\section*{Limitations}

Four caveats apply. First, the LM dictionary is a bag-of-words
measure: it does not capture irony, negation, or context, so the
\emph{specificity} construct used here is a coarse proxy for the
named-entity-based measure of \textcite{hope2016}. Second, the
identification is correlational. The within-firm robustness check
(Section~\ref{sec:robustness_firmfe}) preserves the volume effect
but absorbs both LM sub-densities, which is consistent with
risk-vocabulary composition being a quasi-permanent firm-style
trait \parencite{cazier2021} but means H2 and H3 cannot be claimed
to hold in the time series of a given firm. Third, the sample is
restricted to large U.S.\ non-financial firms with continuous XBRL
coverage; results may not generalise to small caps, financials, or
non-U.S.\ filers. Fourth, the boilerplate-disclosure interpretation
of LitDensity follows the \textcite{nelson2007,cazier2021,beatty2019}
literature but is not directly tested with a measure of safe-harbour
text recycling; it is consistent with the data rather than
demonstrated by them.

\section*{Suggested extensions}

Three extensions follow naturally. First, the named-entity
pipeline of \textcite{hope2016} would replace the dictionary-based
specificity proxy with a direct measure of risk-factor specificity
and would provide a stronger test of H2. Second, a forward-looking
statement classifier in the spirit of \textcite{li2010}, possibly
built on a fine-tuned FinBERT-FLS encoder, would let the H3
literature's distinction between forward-looking cautionary
language and backward-looking litigious narrative be tested
separately; this would also extend the seven null FinBERT-tone
results documented in Section~\ref{sec:robustness_failed} to the
functional form (FLS classification) that the FinBERT family is
actually trained for. Third, an explicit text-recycling measure
(year-on-year overlap of LM-Litigious sentences) would test the
boilerplate-disclosure interpretation of the H3 result directly,
rather than inferring it from the cross-sectional pattern of
LitDensity coefficients.
